\chapter{Analytic signals, circular statistics, and cross-frequency coupling}\label{ch:hilbert}
\begin{readergoal}Connect complex phase to the analytic signal and understand exactly what phase-amplitude coupling measures.\end{readergoal}
\section{The analytic signal}
For a real-valued signal $x(t)$, define its Hilbert transform $\mathcal H\{x\}(t)$ by a principal-value convolution with $1/(\pi t)$, under suitable function-space assumptions. The analytic signal is
\[z(t)=x(t)+i\mathcal H\{x\}(t)=a(t)e^{i\phi(t)},\]
where $a(t)=|z(t)|$ is amplitude and $\phi(t)=\arg z(t)$ is phase. In Fourier space, analytic-signal construction suppresses negative frequencies and doubles positive frequencies (with appropriate DC and Nyquist conventions). These operations acquire a clean oscillatory interpretation for a reasonably narrowband component. Phase can be unstable when instantaneous amplitude is very small.

\section{Filtering comes before interpretation}
Often the signal is band-pass filtered to $x_b(t)$ before Hilbert transformation. The resulting $a_b(t)$ and $\phi_b(t)$ describe the \emph{filtered component}, not a unique property of the raw potential. Choice of passband, filter order, padding, and transitions can change the estimates. For nonsinusoidal waves, harmonics introduced by waveform shape can appear as high-frequency power linked to the low-frequency phase \cite{ColeVoytek2018}.

\section{Circular means and concentration}
For phases $\phi_1,\ldots,\phi_n$, set $R=n^{-1}\sum_j e^{i\phi_j}$. The mean direction is $\arg R$ when $R\ne0$, and the mean resultant length $|R|$ measures concentration between zero and one. Classical tests such as Rayleigh's test target nonuniformity in directions under specified assumptions. They do not automatically correct for temporal autocorrelation or repeated units from the same animal \cite{Berens2009}.

\section{Phase locking and phase-amplitude dependence}
The phase-locking value of phase differences is $\mathrm{PLV}=|n^{-1}\sum_j e^{i(\phi_{1j}-\phi_{2j})}|$. It measures concentration of phase differences, not amplitude modulation. A PAC statistic such as the mean vector length $|n^{-1}\sum_j A_{\mathrm{high}}(t_j)e^{i\phi_{\mathrm{low}}(t_j)}|$ quantifies phase dependence of a high-band amplitude; its scale and biases require normalization and surrogates. The modulation index based on divergence of binned amplitude distributions is a different statistic. These quantities should not be conflated.

\section{The interpretation trap}
A significant PAC value need not represent two physiologically distinct oscillators interacting. Common nonstationarity, sharp transients, harmonics, or filtering artifacts can generate cross-frequency relationships \cite{Aru2015}. Good practice checks waveform shape, aperiodic structure, trial boundaries, and meaningful surrogate null models. Early PAC observations in human cortex motivate the question, but they do not eliminate the general identifiability issues \cite{Canolty2006}.

\begin{workedexample}[One nonsinusoidal oscillator]
Consider the following waveform:
\[x(t)=\cos(2\pi f_0 t)+0.5\cos(4\pi f_0 t).\]
Its higher-frequency harmonic is phase-determined by the same fundamental oscillator. Filtering separately near $f_0$ and $2f_0$ produces components whose phase relationship is stable. It would be misleading to infer two interacting independent biological oscillators from this constructed example.
\end{workedexample}
% --- expanded edition (October 2026) ---
\section{The Hilbert transform in the frequency domain}
With the convention $\hat x(f)=\int x(t)e^{-2\pi ift}dt$, the Hilbert transform acts as the multiplier
\[
\widehat{\mathcal Hx}(f)=-i\operatorname{sgn}(f)\,\hat x(f).
\]
It shifts every positive-frequency component by $-\pi/2$ and every negative-frequency component by $+\pi/2$, leaving magnitudes unchanged. Hence $\mathcal H\cos=\sin$ and $\mathcal H\sin=-\cos$, the analytic signal of $\cos(2\pi f_0t)$ is $e^{2\pi if_0t}$, and $\widehat z(f)=2\hat x(f)\mathbf 1_{f>0}+\hat x(0)\mathbf 1_{f=0}$. Because the multiplier has modulus one, $\mathcal H$ preserves energy and $\mathcal H(\mathcal Hx)=-x$ for signals without a DC component.

\section{When is amplitude-phase decomposition meaningful?}
Any real signal has an analytic representation $a(t)e^{i\phi(t)}$, so the decomposition always exists. The question is whether $a$ and $\phi$ have the intended interpretation as ``envelope'' and ``phase of a rhythm.'' Bedrosian's theorem gives a sufficient condition: if $a(t)$ has spectrum in $[-B,B]$ and the carrier frequency exceeds $B$, then the analytic signal of $a(t)\cos(2\pi f_0t+\theta)$ is $a(t)e^{i(2\pi f_0t+\theta)}$, so the envelope is recovered exactly. In words: the amplitude must vary slowly compared with the oscillation. A theta-band filter of width 4\,Hz centred at 8\,Hz approximately satisfies this; a 1--20\,Hz ``broadband'' filter does not, and its ``phase'' mixes several rhythms.

The \emph{instantaneous frequency} $\frac1{2\pi}\phi'(t)$ is a derivative of phase and inherits noise amplification. When $a(t)$ passes near zero the phase can jump by nearly $\pi$, producing large positive or negative instantaneous frequencies that have no physiological meaning. Such points should be masked using an amplitude threshold.

\section{Distributions on the circle}
The \emph{von Mises} distribution with mean direction $\mu$ and concentration $\kappa\ge0$ has density
\[
p(\theta)=\frac{e^{\kappa\cos(\theta-\mu)}}{2\pi I_0(\kappa)},
\]
where $I_0$ is a modified Bessel function. It plays the role of the normal distribution on the circle: $\kappa=0$ is uniform, and large $\kappa$ approximates $\mathcal N(\mu,1/\kappa)$. Its expected resultant length is $A(\kappa)=I_1(\kappa)/I_0(\kappa)$.

\begin{proposition}[Bias of the resultant length]
If $\phi_1,\ldots,\phi_n$ are iid uniform on the circle, then $\E\bigl[|R|^2\bigr]=1/n$ for $R=\frac1n\sum_je^{i\phi_j}$.
\end{proposition}
\begin{proof}
$|R|^2=\frac1{n^2}\sum_{j,k}e^{i(\phi_j-\phi_k)}$. The $n$ diagonal terms equal $1$; for $j\ne k$, $\E[e^{i(\phi_j-\phi_k)}]=\E[e^{i\phi_j}]\E[e^{-i\phi_k}]=0$.
\end{proof}
Hence $|R|\approx1/\sqrt n$ for pure noise: 25 spikes give an expected PLV near $0.2$ without any locking. Comparisons of PLV across conditions with different spike counts are biased. The \emph{pairwise phase consistency} $\mathrm{PPC}=\frac{2}{n(n-1)}\sum_{j<k}\cos(\phi_j-\phi_k)$ removes the diagonal terms and is an unbiased estimator of the squared population resultant length, $\mathrm{PPC}=(n|R|^2-1)/(n-1)$. The Rayleigh test uses $Z=n|R|^2$, which is approximately exponential with mean one under uniformity, so $P(Z>z)\approx e^{-z}$ for large $n$.

\section{Two PAC estimators and their assumptions}
Let $\phi(t)$ be a low-frequency phase and $A(t)$ a high-frequency amplitude.
\begin{enumerate}[leftmargin=2em]
\item The \emph{mean vector length} $\mathrm{MVL}=|\frac1T\sum_tA(t)e^{i\phi(t)}|$ depends on the overall scale of $A$; it is usually normalized by surrogates.
\item The \emph{modulation index} bins $\phi$ into $J$ bins, computes the mean amplitude $\bar A_j$ in each, normalizes to a distribution $p_j=\bar A_j/\sum_k\bar A_k$, and reports $\mathrm{MI}=D_{\mathrm{KL}}(p\,\|\,u)/\log J$, where $D_{\mathrm{KL}}(p\,\|\,u)=\log J-H(p)$ is the KL divergence from the uniform distribution. It is invariant to scaling $A$ but depends on $J$ and on the number of cycles.
\end{enumerate}
Both estimators are positive under the null and grow with recording length differently, so they must be compared with surrogates that preserve the autocorrelation of each signal (e.g.\ cutting $A$ at a random point and swapping the two pieces), not with zero.

\begin{workedexample}[Envelope of two close sinusoids]
Let $x(t)=\cos(2\pi\cdot 8t)+\cos(2\pi\cdot 9t)=2\cos(\pi t)\cos(2\pi\cdot8.5t)$. Its analytic signal is $e^{2\pi i8t}+e^{2\pi i9t}$, with amplitude $|2\cos(\pi t)|$, a 1\,Hz beat, and phase advancing at 8.5\,Hz but jumping by $\pi$ every second when the amplitude passes through zero. An envelope fluctuation and a ``phase reset'' appear although two perfectly stationary rhythms are present. Amplitude modulation and phase discontinuities in a band-limited signal are therefore not, by themselves, evidence of transient events.
\end{workedexample}
\begin{exercises}
\item Show why the arithmetic mean of angles near $\pm\pi$ is unreliable.
\item Distinguish PLV from PAC with explicit input variables.
\item Propose a surrogate method and specify exactly what structure it preserves.
\item Explain why zero-phase filters cannot be interpreted as real-time causal processing.
\item Show that $\mathcal H\{\cos(2\pi f_0t)\}=\sin(2\pi f_0t)$ for $f_0>0$ using the frequency-domain multiplier.
\item Compute the expected PLV for 100 uniformly distributed spike phases, and compare with the value $0.2$ for 25 phases.
\item Prove the identity $\mathrm{PPC}=(n|R|^2-1)/(n-1)$.
\item Show that the modulation index is zero exactly when the binned amplitude distribution is uniform.
\item In the beat example, find the times at which the phase jumps and the size of the jump.
\end{exercises}
